\begin{ledgroupsized}[r]{120mm}
\footnotesize
\pstart
\noindent\textbf{Überlieferung:}
\pend
\end{ledgroupsized}
\begin{ledgroupsized}[r]{114mm}
\footnotesize
\pstart
\parindent -6mm
\makebox[6mm][l]{\textit{L}}%
Auszüge mit Bemerkungen aus 
\indexN{\textls{Spoleti}, Francesco, gest. 1712}\textsc{F.~Spoleti},
\cite{02095}\textit{De momento quo gravia nituntur deorsum per datum planum} \lbrack...\rbrack\ \textit{de secretione bilis in Hepate}, Venedig 1686:
LH~III~5~Bl.~17\textendash18. 
Ein Bogen~4\textsuperscript{o};
\href{\Scan{LH_03_05_017-018wz}}%
{Wasserzeichen} im Falz;	
oberer Rand beschnitten.
Eine Seite auf Bl.~18~r; Bl.~18~v leer;
Bl.~17 überliefert einen Text mit der Überschrift \glqq De secretione animali\grqq: Dabei handelt es sich um umfangreiche Auszüge mit Bemerkungen aus \indexN{\textls{Spoleti}, Francesco, gest. 1712}Spoletis medizinischer Dissertation (\textit{De secretione bilis in hepate}), die vsl. in einem späteren Band der Reihe erscheinen werden.
\pend
\end{ledgroupsized}
\selectlanguage{latin}
\vspace{8mm}
\pstart
\noindent
\edtext{\href{\Scan{LH_03_05_018r}}%
{\lbrack18~r\textsuperscript{o}\rbrack} 
Francisci}{%
\lemma{\href{\Scan{LH_03_05_018r}}%
{\lbrack18~r\textsuperscript{o}\rbrack}}%
\Bfootnote{%
\textit{(1)}~Ex Franci
\textit{(2)}~Francisci~\textit{L}}}
\indexN{\textls{Spoleti}, Francesco, gest. 1712}Spoleti 
\edtext{\cite{02095}\textit{De Momento quo gravia nituntur deorsum}
\edtext{\textit{per datum planum,}}{\lemma{}\Bfootnote{\textit{per datum planum,} \textit{erg.~L}}}
\textit{et quo subjectum planum urgent, item quo planum inclinatum grave deorsum conatur.
Altera ejusdem Dissertatio de secretione bilis in Hepate}
Venet. 1686. 4\textsuperscript{o} apud Andr. Poleti.}{%
\lemma{\textit{De momento} \lbrack...\rbrack\ Poleti}%
\Cfootnote{\indexN{\textls{Spoleti}, Francesco, gest. 1712}\textsc{F.~Spoleti},
\cite{02095}\textit{De momento quo gravia nituntur deorsum per datum planum} \lbrack...\rbrack\ \textit{de secretione bilis in Hepate}, Venedig 1686, Titelseite.}}
\pend
\pstart
\edtext{\indexN{\textls{Christina}, Königin von Schweden 1632\textendash1654, gest. 1689}Christinae reginae dedicat.}{%
\lemma{}%
\Bfootnote{Christinae reginae dedicat. \textit{erg.~L}}%
\lemma{Christinae reginae dedicat}%
\Cfootnote{\cite{02095}a.\,a.\,O., S.\;3\,f.}}
\pend
\pstart
\edtext{Ait se anno superiore \indexO{Rom}Romae edidisse
\textit{Unicam propositionem, in qua institutum ipsi fuerit ostendere quod 
Momentum gravis in plano inclinato\indexS{momentum gravis in plano inclinato} 
esse ad momentum totale,\indexS{momentum gravis totale} 
ut perp.\indexS{perpendiculum plani inclinati} ad Hypoth.\indexS{hypothenusa plani inclinati}}}{%
\lemma{Ait \lbrack...\rbrack\ \textit{Hypoth.}}%
\Cfootnote{\cite{02095}a.\,a.\,O., S.\;5 mit Änderungen; siehe \indexN{\textls{Spoleti}, Francesco, gest. 1712}\textsc{Spoleti}, \cite{02096}\textit{De momento gravis in plano inclinato}, Rom 1684.}}
\edtext{%
\textit{Demonstrationem a nonnullis damnatam obscuritatis} qui 
\textit{\indexN{\textls{Archimedes} 287\textendash212 v. Chr.}Archimedis et 
\indexN{\textls{Valerio}, Luca 1552\textendash1618}Lucae Valerii libros non bene calleant}, 
a quibus dum \textit{Pseudocatonibus qui papyraceam in \indexO{Rom}urbe Mathesin jactant}
pro \textit{spuria} habitam, a nonnullis 
\edtext{culpatam, quasi}{%
\lemma{culpatam,}%
\Bfootnote{%
\textit{(1)}~quod
\textit{(2)}~quasi~\textit{L}}}
trigonometriam\indexS{trigonometria} et quadraturam Circuli\indexS{quadratura circuli} supponeret.
Hic verum \textit{edit cum lemmatibus} simulque ait 
\textit{novam} a se exhiberi 
\textit{methodum}
\edtext{\lbrack\textit{inveniendi}\rbrack}{%
\lemma{}%
\Bfootnote{%
exhibendi
\textit{L ändert Hrsg. nach Vorlage}}}
\textit{centra Gravitatis\indexS{centrum gravitatis segmenti sphaerae} datorum Spherae segmentorum.}}{%
\lemma{\textit{Demonstrationem} \lbrack...\rbrack\ \textit{segmentorum}}%
\Cfootnote{\indexN{\textls{Spoleti}, Francesco, gest. 1712}\textsc{F.~Spoleti},
\cite{02095}\textit{De momento} \lbrack...\rbrack\ \textit{de secretione bilis}, S.\;5, mit Änderungen.
Siehe
\indexN{\textls{Archimedes} 287\textendash212 v. Chr.}\textsc{Archimedes}, 
\cite{01010}\textit{De planorum aequilibris libri duo};
\indexN{\textls{Valerio}, Luca 1552\textendash1618}\textsc{L.~Valerio}, 
\cite{02100}\textit{De centro gravitatis solidorum libri tres}, Rom 1604.}}
\edtext{Ait se eodem usurum Schemate quo 
\indexN{\textls{Vanni}, Giovanni Francesco 1638\textendash1709}doctor Lucensis
qui receptam de plano inclinato\indexS{planum inclinatum} sententiam\indexS{sententia recepta de plano inclinato} impugnare ausus est
in suo \cite{01995}\textit{Specimine} et \cite{02000}\textit{Exegesi Physico-Mathematica} 
est usus.}{%
\lemma{Ait \lbrack...\rbrack\ usus}%
\Cfootnote{\cite{02095}a.\,a.\,O., S.\;7.
Vgl. \indexN{\textls{Vanni}, Giovanni Francesco 1638\textendash1709}\textsc{G.\,F.~Vanni},
\cite{01995}\textit{Specimen libri de momentis gravium}, \textit{AE} November 1684, S.\;511\textendash513
und \cite{02000}\textit{Exegeses physico-mathematicae}, Rom 1685.
Zu Leibnizens Reaktionen auf diese Publikationen Vannis siehe u.\,a. die (anonym) veröffentlichten Stücke N.\;\ref{intro_vanni-specimen}, N.\;\ref{49409} und N.\;\ref{41199}. 
Das hier erwähnte Diagramm, das \indexN{\textls{Spoleti}, Francesco, gest. 1712}Spoleti dem \textit{Specimen} (dort auf S.\;512) entnimmt, wird auch von Leibniz in seinen Besprechungen von \indexN{\textls{Vanni}, Giovanni Francesco 1638\textendash1709}Vanni wiederaufgegriffen.}}
\pend
\begin{edfigure}%
\includegraphics[scale=1]{%
gesamttex/edit_VIII-4/images/LH_03_05_018_d.pdf}
\edcaption[0em]{\edtext{\lbrack\textit{Fig.~1}\rbrack}{%
\lemma{\lbrack\textit{Fig.~1}\rbrack}%
\Cfootnote[nonum]{Ein gestrichener Entwurf zum Diagramm wird nicht wiedergegeben.}}}
\end{edfigure}
\pstart
\edtext{Ostendit sphaeram\indexS{sphaera} in plano inclinato\indexS{planum inclinatum} librari super puncto 
\edtext{insistentiae\indexS{punctum insistentiae} \textit{E}, plano \textit{EC} recto ad horizontem\lbrack,\rbrack\ 
et planum inclinatum\indexS{planum inclinatum} secari in duo segmenta, \textit{ACE}, et \textit{FCE}}{%
\lemma{insistentiae \textit{E},}%
\Bfootnote{%
\textit{(1)}~et plano verticali ad horizontem et recto ad planum
\textit{(2)}~\bars{et \textit{streicht Hrsg.}} plano
\textit{(a)}~ver
\textit{(b)}~\bars{\textit{EC} \textit{erg.}} recto \lbrack...\rbrack\ \textit{FCE}~\textit{L}}} 
et volvi deorsum conari 
\edtext{excessu, quo}{%
\lemma{excessu}%
\Bfootnote{%
\textit{(1)}~momenti, quod
\textit{(2)}~, quo~\textit{L}}}
momentum\indexS{momentum} segmenti
\edtext{majoris}{\lemma{}\Bfootnote{majoris \textit{erg.~L}}}
\textit{FCE} excedit momentum contrarium
\edtext{segmenti}{\lemma{}\Bfootnote{segmenti \textit{erg.~L}}}
\textit{ACE}.
Calculo autem ex doctrina de centris gravitatis segmentorum sphaerae\indexS{centrum gravitatis segmenti sphaerae} ducto
ostendit differentiam horum ductum momentorum esse ad momentum totale sphaerae\indexS{momentum sphaerae totale}
\edtext{ut perpendiculum}{%
\lemma{ut}%
\Bfootnote{%
\textit{(1)}~cathetum
\textit{(2)}~perpendiculum~\textit{L}}}
plani inclinati\indexS{perpendiculum plani inclinati} ad 
\edtext{Hypothenusam,\indexS{hypothenusa plani inclinati} seu}{%
\lemma{Hypothenusam,}%
\Bfootnote{%
\textit{(1)}~seu ut \textit{MB} ad \textit{C}
\textit{(2)}~nam
\textit{(3)}~seu~\textit{L}}}
ut \textit{MB} ad \textit{MA}
\edtext{seu ad \textit{ME}}{\lemma{}\Bfootnote{seu ad \textit{ME} \textit{erg.~L}}}
nam, quo magis erigitur planum \textit{PN}
eo magis \textit{B} accedit ad \textit{A}, donec coincidunt cum planum est rectum \textit{QA}.
Sunt autem triangula \textit{MBE} et \textit{PRE} similia.}{%
\lemma{Ostendit \lbrack...\rbrack\ similia}%
\Cfootnote{\indexN{\textls{Spoleti}, Francesco, gest. 1712}\textsc{Spoleti},
\cite{02095}\textit{De momento gravis}, S.\;17\,f.}}
\pend
\pstart
\edtext{%
Idem et brevius sic ostendit\lbrack:\rbrack\ Momentum sphaerae in plano inclinato,\indexS{momentum sphaerae in plano inclinato} 
esse aestimandum distantia ejus centri gravitatis\indexS{centrum gravitatis} \textit{M},
\edtext{a verticali}{\lemma{}\Bfootnote{a verticali \textit{erg.~L}}}
per punctum librationis\indexS{punctum librationis} \textit{E} transeunte \textit{EC}\lbrack;\rbrack\ 
quae distantia est \textit{MB} et plano erecto\indexS{planum erectum} abit in radium.%
}{\lemma{Idem \lbrack...\rbrack\ radium}%
\Cfootnote{\cite{02095}\textit{a.\,a.\,O.}, S.\;19.}}
\pend
\pstart
\edtext{%
Notat praeterea momentum quo sphaera descendit in plano,\indexS{momentum quo sphaera descendit in plano} additum momento quo
\edtext{sphaera premit}{%
\lemma{sphaera}%
\Bfootnote{%
\textit{(1)}~trudit
\textit{(2)}~in pla
\textit{(3)}~premit~\textit{L}}}
ipsum planum,\indexS{momentum quo sphaera premit planum} aequari totali momento sphaerae.\indexS{momentum sphaerae totale}%
}{\lemma{Notat \lbrack...\rbrack\ sphaerae}%
\Cfootnote{\cite{02095}a.\,a.\,O., S.\;19\,f.}}
\edtext{%
Notat et rectam inclinatam et circa centrum quo subjecto plano insistit mobilem\indexS{recta inclinata circa centrum mobilis}
momentum quo descendere conatur trahere, quod sit similiter
\edtext{ad momentum ejus totale}{\lemma{}\Bfootnote{ad momentum ejus totale \textit{erg.~L}}}
ut perpendicularis inclinationis\indexS{perpendicularis inclinationis} ad Hypothenusam.\indexS{hypothenusa inclinationis}%
}{\lemma{Notat \lbrack...\rbrack\ Hypothenusam}%
\Cfootnote{\cite{02095}a.\,a.\,O., S.\;21\,f.}}
\edtext{Utileque id esse ad motus \textit{pontium versatilium\indexS{pontes versatiles} qui ad arcium custodiam} adhibentur aestimandos.}{%
\lemma{Utileque \lbrack...\rbrack\ aestimandos}%
\Cfootnote{\cite{02095}a.\,a.\,O., coroll. \lbrack II\rbrack, S.\;22 mit Änderungen.}}
\pend
\pstart
Dissertatio ejusdem \textit{de secretione Bilis\indexS{secretio bilis} in Hepate\indexS{hepar}}
brevis quidem est, sed quaedam habet peringeniosa.
\edtext{Notat 
\indexN{\textls{Cesalpino} (Caesalpinus), Andrea 1524 od. 1525\textendash1603}\textit{Caesalpinum
detexisse sanguinem\indexS{sanguis} a dextro\indexS{ventriculus dexter} per pulmones\indexS{pulmo} in sinistrum\indexS{ventriculus sinistrus} cordis\indexS{cor} ventriculum ruentem}
idque \indexN{\textls{Harvey} (Harvaeus), William 1578\textendash1657}Harvaeo lucem attulisse ut 
\textit{detegeret sanguinem\indexS{sanguis} per aortam\indexS{aorta} et reliquas arterias\indexS{arteria} in universum corpus
impelli}\lbrack,\rbrack\ \textit{per venas\indexS{vena} iterum} ad cor\indexS{cor} \textit{regressurum}.}{%
\lemma{Notat \lbrack...\rbrack\ regressurum}%
\Cfootnote{\cite{02095}a.\,a.\,O., S.\;23 mit Auslassungen.
Siehe \indexN{\textls{Cesalpino} (Caesalpinus), Andrea 1524 od. 1525\textendash1603}\textsc{A.~Cesalpino}, \cite{02111}\textit{Quaestionem medicarum libri II}, Venedig 1593;
\indexN{\textls{Harvey} (Harvaeus), William 1578\textendash1657}\textsc{W.~Harvey},
\cite{02112}\textit{Exercitatio anatomica de motu cordis et sanguinis in animalibus}, Frankfurt 1628.}}
\edtext{Borelli\indexN{\textls{Borelli} (Borellus), Giovanni Alfonso 1608\textendash1679} opus 
\cite{02099}\textit{De motu animalium} magnifacit\lbrack;\rbrack\ porro Mechanicam\indexS{Mechanica}
\edlabel{LH_03_05_018_1a}%
\edtext{}{%
\xxref{LH_03_05_018_1a}{LH_03_05_018_1b}%
\lemma{et}%
\Bfootnote{%
\textit{(1)}~chymicam
\textit{(2)}~chemicam esse fungendas.
\textit{(a)}~Sed
\textit{(b)}~Non~\textit{L}}}%
et chemicam\indexS{chemica} esse fungendas.}{%
\lemma{Borelli \lbrack...\rbrack\ fungendas}%
\Cfootnote{%
\indexN{\textls{Spoleti}, Francesco, gest. 1712}\textsc{Spoleti},
\cite{02095}\textit{De momento} \lbrack...\rbrack\ \textit{de secretione bilis}, S.\;23\,f.
Siehe \indexN{\textls{Borelli} (Borellus), Giovanni Alfonso 1608\textendash1679}\textsc{G.\,A.~Borelli},
\cite{02099}\textit{De motu animalium}, Rom 1680\textendash1681.}}
\edtext{\textit{Non\edlabel{LH_03_05_018_1b} dari quatuor humores,\indexS{humor} sed unum ex pluribus constitutum, 
gignere in glandulosa cerebri\indexS{cerebrum} substantia\indexS{substantia glandulosa cerebri}
nerveum succum\indexS{succus nerveus} cui adscribendae facultatis animalis operationes,\indexS{operationes facultatis animalis} gignere in ventriculo\indexS{ventriculus} liquidum
nitro-sulphureum\indexS{liquidum nitro-sulphureum} quo diluitur cibus\indexS{cibus}} et digeritur
non \textit{calore innato},\indexS{calor innatus} qui \textit{figmentum} sit, quemadmodum \textit{et humidum radicale}.\indexS{humidum radicale}}{%
\lemma{\textit{Non dari} \lbrack...\rbrack\ \textit{radicale}}%
\Cfootnote{%
\indexN{\textls{Spoleti}, Francesco, gest. 1712}\textsc{Spoleti},
\cite{02095}\textit{De momento} \lbrack...\rbrack\ \textit{de secretione bilis}, S.\;24. mit Auslassungen.}}
\edtext{Cum alia \textit{ex arterioso sanguine}\indexS{sanguis arteriosus} separentur,
\textit{sola bilis\indexS{bilis} secernitur ex} 
\edtext{\lbrack\textit{venoso}\rbrack.\indexS{sanguis venosus}}{%
\lemma{}%
\Bfootnote{seroso \textit{L ändert Hrsg. nach Vorlage}}}%
}{%
\lemma{Cum alia \lbrack...\rbrack\ \lbrack\textit{venoso}\rbrack}%
\Cfootnote{\cite{02095}a.\,a.\,O., S.\;24 mit Auslassungen.}}
\edtext{Non \textit{fieri secretiones\indexS{secretio} per solam cribrationem,\indexS{cribratio}}
quasi fluida sine ullo \textit{nexu et solo contactu simul haereant,} 
et pororum ratione qua diversis partibus congruant discernantur;}{%
\lemma{Non \textit{fieri} \lbrack...\rbrack\ discernantur}%
\Cfootnote{\cite{02095}a.\,a.\,O., S.\;24\,f. mit Auslassungen.}}
\edtext{nec fermentatione\indexS{fermentatio}
\edtext{sola}{\lemma{}\Bfootnote{sola \textit{erg.~L}}}
fieri ut chymici\indexS{chymici} 
\edtext{volunt\lbrack,\rbrack\ quomodo}{%
\lemma{volunt}%
\Bfootnote{%
\textit{(1)}~unde
\textit{(2)}~quomodo~\textit{L}}}
enim ipsum fermentum prius secretum?
Nec \textit{motum illum intestinum vertiginosum} fermentationi\indexS{fermentatio}
necessarium \textit{fieri} omnimode \textit{posse in angustissimis illis canaliculis.}%
}{%
\lemma{nec fermentatione \lbrack...\rbrack\ \textit{canaliculis}}%
\Cfootnote{\cite{02095}a.\,a.\,O., S.\;25\,f. mit Auslassungen.}}
\pend
